\subsection*{Note (7)}
The final note, on pp.~56--57, supplies the data behind the flag-indexed
system of spaces mentioned in \S5, p.~34. Its ``strates déployées'',
``tubes de raccord'', ``fibres de raccord'', and ``espaces de jonction''
describe the unfolding of the stratified structure. A target-index slip
occurs in the prose on p.~56: the first morphism is written
$X_{i,j}\longrightarrow X_i$ and described as a bordering inclusion;
the second is written $X_{i,j}\rightarrow X_i$ and described as a proper
fibration. The diagram and its stated legend
instead give the hooked immersion
$a_{ij}:X_{ij}\hookrightarrow X_j$ and the ordinary proper fibration
$b_{ij}:X_{ij}\longrightarrow X_i$. Thus it is the target of the first
prose morphism that conflicts with the diagram; both printed prose
targets remain in the readers. The lower right cartesian square identifies
$X_{ijk}$ with $X_{ij}\times_{X_j}X_{jk}$ through its two projections,
using $a_{ij}$ and $b_{jk}$ as the maps to $X_j$. The author's NB stresses
that the data still include the morphism
$X_{ijk}\longrightarrow X_{ik}$, since the immersion into the target
$X_k$ is not required to be an embedding. The scan has two faint marks
forming a colon-like trace at subscript height after this final $k$; their
physical origin and textual role remain unresolved. On p.~57
he recasts the unfolding of a stratified tame space using the sum $X_*$
of its unfolded strata and the order relation with graph $X_{**}$:
$pr_1:X_{**}\longrightarrow X_*$ is the proper locally trivial fibration,
and $pr_2:X_{**}\longrightarrow X_*$ the bordering embedding. For tame
multiplicities the account uses a category structure in place of that
ordered-object structure, with the maps
$X_{ijk}\longrightarrow X_{ik}$ furnishing composition. The note thus
turns the extra structures and reconstruction problem of \S5 into
explicit morphism data and clarifies the role of the junction spaces.
(\href{https://grothendieck.umontpellier.fr/web/viewer.html?file=../119.pdf\#page=59}{Copy 1, pp.~56--57, beginning at site 58};
\href{https://grothendieck.umontpellier.fr/web/viewer.html?file=../119.pdf\#page=60}{ordered and categorical descriptions on p.~57, site 59};
\href{https://grothendieck.umontpellier.fr/web/viewer.html?file=../119.pdf\#page=37}{call in \S5, p.~34, site 36}.)
